% Shared body for Paper B: "SoK: Authorization for Autonomous Agents".
% Class-agnostic: the per-venue main.tex sets class/title/authors and provides the
% bibliography. IEEEtran/llncs place \begin{abstract} after \maketitle, which is how
% the wrappers input this file. Unfinished empirical results are wrapped in \gate{}.
\ifdefined\ifsubmission\else\expandafter\newif\csname ifsubmission\endcsname\submissiontrue\fi
\providecommand{\gate}[1]{\ifsubmission #1\else \textbf{[PENDING: #1]}\fi}
\providecommand{\draftonly}[1]{\ifsubmission\else \textbf{[NOTE: #1]}\fi}
\providecommand{\yes}{\ensuremath{\surd}}
\providecommand{\no}{\ensuremath{\times}}
\emergencystretch=3em

% Abstract in abstract.tex; acmart needs it before \maketitle (emitted by the wrapper),
% so the body only emits it when the wrapper has not.
\ifdefined\keywayabstractdone\else\begin{abstract}
An autonomous agent calls tools and Model Context Protocol (MCP) servers while
carrying a bearer token (an OAuth~2.0 access token or an MCP authorization token)
on each request. Because the agent is driven by untrusted input, its authorization
layer is part of the attack surface: a token not bound to the server it is sent to
can be replayed or passed through, an over-broad scope hands an injected agent far
more than any task needs, and an on-behalf-of token with no verifiable delegation
chain leaves ``who authorized this?'' unanswerable. A fast-growing research
literature studies agent security, but its coverage is densest on model-level attacks
(prompt injection, tool poisoning) and it proposes individual mechanisms; it does not
systematize the authorization layer. We survey 44 research works and show that two
authorization properties, audience binding and confused-deputy / dynamic client
registration, are treated as mere specification obligations rather than as
properties anyone tests or enforces. We then systematize agent authorization along
three axes no prior work combines: (i) by \emph{pre-deployment testability}, meaning
which threats are decidable from a token, a key set, and a declared tool contract
before the agent runs; (ii) by \emph{normative provenance}, a threat-to-clause mapping to
the OAuth RFCs, the MCP authorization specification, OWASP, and CWE; and (iii) by a
\emph{three-way enforcement spectrum}, token-static vs.\ server-dynamic vs.\ needing a
human or the model. Across fifteen threats in six categories, we find twelve are
testable before production, six from a token alone and six against a staging server
with no human or model in the loop, and only three need a human or the model. We
evaluate the token-static subset against nine independent verifiers, eight libraries
across three ecosystems and the Envoy proxy, and find six accept a token they should
reject, a no-expiry token (four, including the Envoy proxy that fronts service meshes)
or an unbound-audience token (one), so the weaknesses are defaults, not corner cases. We argue this cut should guide both tooling and standardization.
\end{abstract}
\fi

\section{Introduction}
When an agent calls a tool or an MCP server on a user's behalf, a bearer token
travels with the request. The failure modes are OAuth failure modes, but with
sharper consequences, because the caller is autonomous and can be steered by
untrusted input such as a retrieved document or a tool description~\cite{agentauth2023indirectpi}.
A token not bound to the resource it is presented to can be replayed against, or
forwarded through, a different server (a confused deputy); an on-behalf-of (OBO)
token with no verifiable delegation chain makes accountability impossible; an
omnibus scope such as \texttt{admin:*} hands an injected agent the keys to
everything; and a non-expiring credential never stops being useful to whoever
obtains it.

\noindent\textbf{What agent autonomy and MCP change.} Each failure above is a known
OAuth or JWT weakness, and the individual checks that catch them, require an
\texttt{aud}, cap a scope, demand an \texttt{exp}, are classical hygiene. Two things
are new. First, the client is now autonomous and driven by untrusted input, so a
single over-broad or unbound token is not merely a misconfiguration a careful client
avoids but a live capability an injected plan will exercise: the token is the attack
surface. Second, MCP composes what single-service OAuth did not, and the composition
is where authorization breaks: tokens pass through chains of servers, delegation nests
across agent and sub-agent, clients register dynamically, and a server publishes a
declared tool contract against which a token's scope can finally be judged. Our
contribution is therefore not the individual checks but the systematization they
enable under this threat model, sorting the threats by who can test each, and before
what.

A large and fast-moving body of research now studies agent and MCP security. As our
survey (\S\ref{sec:related}) shows, it is dense on \emph{model-level} attacks
(prompt-injection-driven escalation and tool poisoning) and it proposes many
\emph{individual} defense mechanisms (privilege controllers, delegation protocols,
identity architectures). What it does not do is systematize the \emph{authorization}
layer itself: no surveyed work sorts agent-authorization threats by whether a
defender can test them before deployment, maps each threat to the specific normative
clause it violates, or names the line between what can be enforced statically at the
trust boundary and what can only be enforced at runtime. The two authorization
properties that matter most for token safety, audience binding and the
confused-deputy / dynamic-client-registration pair, are treated across 44 works
as specification obligations to cite rather than as properties to verify
(\S\ref{sec:related}).

\noindent\textbf{Contributions.}
(i) We give a source-grounded threat model of fifteen agent-authorization threats in
six categories, each with its broken invariant, the conditions under which it is
exploitable, a mitigation, and its normative source (\S\ref{sec:background},
\S\ref{sec:tax}).
(ii) We systematize this space along three axes no prior work combines: testability,
normative provenance, and a three-way enforcement spectrum, token-static vs.
server-dynamic vs. human-or-model; by it, twelve of fifteen threats are testable
before production and only three need a human or the model (\S\ref{sec:syst}).
(iii) We survey 44 research works on agent authorization and make the coverage gap
precise: audience binding and confused-deputy are the under-formalized seam
(\S\ref{sec:related}).
(iv) We argue which authorization properties a standard or a pre-deployment tool
could mandate, release an open-source analyzer for the token-static subset
(\S\ref{sec:syst}), and evaluate it against nine reference verifiers, six of which
accept a token they should reject (\S\ref{sec:eval}).

\section{Background}\label{sec:background}
\noindent\textbf{Tokens and claims.} An OAuth~2.0 access token authorizes a client to
call a resource on a user's behalf; in agent deployments it is usually a JSON Web
Token (JWT), a signed set of claims. The claims that govern authorization are the
issuer (\texttt{iss}), the intended audience (\texttt{aud}, the resource the token is
\emph{for}), the granted \texttt{scope}, the expiry (\texttt{exp}), and, for
delegation, the authorized party (\texttt{azp}, naming the client) and the actor
(\texttt{act}, naming who is acting). A verifier must check the signature against the
issuer's keys, that \texttt{aud} names itself, that the algorithm is on an
allow-list, and that required claims hold (RFC~8725). RFC~9068 profiles JWT access
tokens and makes audience and issuer validation mandatory for that profile; RFC~8707
(resource indicators) lets a client request a token bound to a specific resource.

\noindent\textbf{Delegation and the \texttt{act}/\texttt{azp} claims.} OAuth token
exchange (RFC~8693) issues a delegated token carrying an \texttt{act} claim that
names the acting party and a \texttt{may\_act} claim stating who is \emph{permitted}
to act for the subject. A single \texttt{act} claim records one hop; a chain
(user~$\rightarrow$~agent~$\rightarrow$~sub-agent) requires \emph{nested}
\texttt{act} claims, and transitive delegation is only verifiable if each exchange
enforces \texttt{may\_act} and the verifier walks the full nested chain. \texttt{azp}
names the client that obtained the token but does not by itself establish a
delegation chain: it identifies the immediate party, not the sequence of
authorizations behind it. This is why \texttt{azp}/\texttt{act} alone are necessary
but not sufficient: without enforced \texttt{may\_act} and nested-chain validation,
``who authorized this?'' remains unanswerable across more than one hop.

\noindent\textbf{MCP, identity, and the confused deputy.} The Model Context Protocol
standardizes how an agent discovers and calls tools; its authorization specification
inherits OAuth and \emph{prohibits} passing a client's token through to an upstream
API. Dynamic client registration (DCR, RFC~7591) lets clients register at runtime,
which an MCP proxy may expose. Workload-identity substrates (WIMSE; SPIFFE/SPIRE)
would give each agent instance a verifiable, short-lived identity distinct from the
human's, but are rarely deployed. The confused-deputy problem (CWE-441), where a more
privileged party misuses its authority on behalf of a less privileged one, is the
classical frame for token passthrough and consent-flow abuse.

\noindent\textbf{The agent trust model.} Two properties of the agent setting sharpen
every OAuth failure above. First, the caller is \emph{autonomous}: it chooses which
tool to call and with which arguments, so a credential it holds is exercised without a
human reviewing each use. Second, its inputs are \emph{untrusted}: an agent reads
documents, web pages, and tool descriptions an attacker may control, so prompt
injection can steer it into requests the user never intended~\cite{agentauth2023indirectpi}.
Together these mean the agent must be treated as a party that can be induced to misuse
whatever its token permits, the assumption our threat model adopts and that recent
multi-agent security agendas make explicit~\cite{agentauth2025multiagentchallenges}.
The consequence for authorization is concrete: the token an agent carries, not a
human's judgement at call time, is the last line of defense, so the properties that
token enforces are what bound the blast radius of an injected agent. This is why we
treat the token-level, statically checkable properties as first-class rather than as
spec footnotes.

\section{A threat model for agent authorization}\label{sec:tax}
We organize the space into six categories. Table~\ref{tab:tax} lists fifteen threats
with the invariant that must hold, a one-line mitigation, and the weakest harness that
can decide the threat before production (its Test class, \S\ref{sec:syst}). We then
give, per category, the exploit condition and the defense.

\noindent\textbf{Attacker and defender capabilities.} Two capabilities must be kept
apart, because conflating them is easy and misleads. The \emph{attacker} we model
steers the agent through untrusted input, a poisoned document or a malicious tool
description, and may capture a token in transit; it does not hold the issuer's signing
key, so it cannot mint a token with an arbitrary \texttt{aud}, strip \texttt{exp}, or
forge an \texttt{act} chain. What it exploits is a token the deployment already
accepts: an over-broad scope it can cash in (SCOPE-01), an unbound or shared-issuer
audience it can replay or pass through (MCP-01/02), a missing delegation record that
defeats attribution (DEL-01). The \emph{defender}, by contrast, red-teams with issuer
control in a test harness: it mints the adversarial token an attacker would benefit
from and checks that its own verifier rejects it. The testability classes of
\S\ref{sec:syst} describe this \emph{defender} capability, not the attacker's: a
threat is token-static when the defender can settle it from a minted token offline,
whether or not a real attacker could have produced that token unaided.

\begin{table*}[t]
\centering\footnotesize
\caption{Agent-authorization threats by category, with the invariant that must hold,
a mitigation, the normative source, and the \emph{weakest harness} that can decide the
threat before production (Test): \textsf{T}, from a token, key set, and declared
contract alone; \textsf{S}, against a staging server or issuer with no human or model
in the loop; \textsf{H}, only with a live human approver or the model itself. Sev.:
severity (critical/high/med, per \S\ref{sec:related}).}
\label{tab:tax}
\begin{tabular}{@{}llp{4.1cm}p{4.3cm}lc@{}}
\toprule
ID & Sev. & Broken invariant (exploit) & Mitigation & Source & Test \\
\midrule
\multicolumn{6}{@{}l}{\emph{Audience binding}}\\
MCP-01 & crit & token accepted with wrong/any audience (passthrough, replay) & verify \texttt{aud} names this resource; reject forwarded tokens & RFC 8707/9728; MCP~\cite{rfc8707,rfc9728,mcpauth} & \textsf{T} \\
MCP-02 & high & token carries no resource indicator (unbound audience) & request resource-bound tokens; require \texttt{aud} & RFC 8707~\cite{rfc8707} & \textsf{T} \\
\multicolumn{6}{@{}l}{\emph{Delegation}}\\
DEL-01 & high & missing or forged \texttt{act} claim (no accountable actor) & require a signed \texttt{act} claim on OBO tokens & RFC 8693~\cite{rfc8693} & \textsf{T} \\
DEL-02 & high & unverifiable nested delegation chain (transitive-trust abuse) & validate the full nested \texttt{act} chain & RFC 8693~\cite{rfc8693} & \textsf{T} \\
DEL-03 & high & \texttt{may\_act} not enforced on exchange & enforce \texttt{may\_act} at the token endpoint & RFC 8693~\cite{rfc8693} & \textsf{S} \\
DEL-04 & med  & impersonation where delegation was required & issue delegated (not impersonated) tokens & RFC 8693~\cite{rfc8693} & \textsf{S} \\
\multicolumn{6}{@{}l}{\emph{Scope}}\\
SCOPE-01 & high & over-scoped agent credential (\texttt{admin:*}) & least-privilege scopes per task & RFC 9700; OWASP~\cite{rfc9700,owaspagentic} & \textsf{T} \\
SCOPE-02 & high & non-expiring / long-lived token & short \texttt{exp}; rotation & RFC 8725~\cite{rfc8725} & \textsf{T} \\
SCOPE-03 & high & agent reuses the human's root credential & distinct agent credential & OWASP~\cite{owaspagentic} & \textsf{S} \\
\multicolumn{6}{@{}l}{\emph{Agent identity}}\\
AID-01 & high & session used as authentication & bind to a verifiable workload identity & WIMSE/SPIFFE~\cite{wimse,spiffe} & \textsf{S} \\
AID-02 & med  & no distinct workload/agent identity & issue per-agent SVID/credential & WIMSE/SPIFFE~\cite{wimse,spiffe} & \textsf{S} \\
\multicolumn{6}{@{}l}{\emph{Confused deputy}}\\
CD-01 & crit & proxy acts on a client's behalf without per-client consent & per-client consent; exact \texttt{redirect\_uri} match & CWE-441; MCP~\cite{cwe441,mcpauth} & \textsf{H} \\
CD-02 & high & dynamic client registration abuse & vet/limit DCR; bind registrations & RFC 7591~\cite{rfc7591} & \textsf{S} \\
\multicolumn{6}{@{}l}{\emph{Agency}}\\
AGT-01 & crit & prompt-injection-driven privilege escalation & per-action authorization; human-in-loop & OWASP~\cite{owaspllm,owaspagentic} & \textsf{H} \\
AGT-02 & high & tool poisoning / rug-pull via tool metadata & vet tool descriptions; pin versions & OWASP~\cite{owaspagentic} & \textsf{H} \\
\bottomrule
\end{tabular}
\end{table*}

\noindent\textbf{Audience binding (MCP-01/02).} A bearer token names, in its
\texttt{aud} claim, the resource it is for. The exploit has two shapes. In
\emph{passthrough} (MCP-01), an MCP server that receives a client's token forwards it
unchanged to an upstream API; because the upstream often shares an issuer it accepts
the token, and the MCP server has acted as a confused deputy wielding the full
authority of the user's token. In \emph{replay} (MCP-02), a token minted with no
resource indicator, or one whose \texttt{aud} a verifier never checks, is accepted at
any server the issuer serves, so an attacker who captures one token reaches many.
Both are properties of claims the token carries and of the server's declared
identity, so a defender mints the adversarial token and confirms rejection before
deployment. The defense is to request resource-bound tokens (RFC~8707), to verify
\texttt{aud} names this server, and to honor the MCP prohibition on passthrough.
Among the works we survey only Agentic JWT~\cite{agentauth2025agenticjwt} and one
enterprise framework~\cite{agentauth2025enterprisemcp} treat audience binding as a
property to enforce rather than a specification clause to cite (\S\ref{sec:related}).

\noindent\textbf{Delegation (DEL-01..04).} When an agent acts for a user, an
on-behalf-of (OBO) token should carry an \texttt{act} claim naming the acting party
(RFC~8693). A single \texttt{act} records one hop; a chain
user~$\rightarrow$~agent~$\rightarrow$~sub-agent needs \emph{nested} \texttt{act}
claims. A caveat governs what any inspection of them can conclude. RFC~8693 defines
the \texttt{act} claim, including its nested prior actors, as \emph{informational}; it
is a record of identity, not an authorization input, and must not on its own drive an
access-control decision. A signed \texttt{act} history therefore shows \emph{who}
acted but does not prove that each hop was \emph{authorized}; that holds only if every
token exchange along the chain enforced the \texttt{may\_act} claim at issuance. The
exploit: with a missing or malformed \texttt{act} (DEL-01) an action cannot be
attributed, and with a chain whose \texttt{may\_act} was never enforced (DEL-02) a
sub-agent silently exceeds what was delegated, a transitive-trust escalation. What a
presented token settles is thus a hygiene property: a verifier confirms from the token
alone that a required \texttt{act} claim is present and the nested chain is
well-formed, and rejects a token that asserts delegation but carries none. Confirming
that the chain was \emph{authorized} is an extension resting on an explicit trust
assumption, that the issuer enforced \texttt{may\_act}, and cannot be read off the
token. The remaining two are plainly runtime: whether the token endpoint enforced
\texttt{may\_act} on each exchange (DEL-03), and whether a delegated rather than an
impersonated token was issued (DEL-04), are properties of the issuing flow, not of the
presented token. Delegation protocols that bind each action to user intent and bound a
sub-agent's authority target this
category~\cite{agentauth2025agenticjwt,agentauth2026aip,agentauth2026boundedagents,agentauth2026aidguard};
an empirical gap analysis finds identity, authorization, and runtime governance thinly
deployed across multi-agent systems in practice~\cite{agentauth2026delegationwithouttrust}.

\noindent\textbf{Scope (SCOPE-01..03).} The \texttt{scope} and \texttt{exp} a token
carries bound what an injected agent can do with it. The exploit: an omnibus scope
such as \texttt{admin:*} (SCOPE-01) hands an injected agent far more than any task
needs, and a non-expiring token (SCOPE-02) stays useful to whoever later obtains it.
Both are read directly from the token, so a check against a declared allow-list and a
maximum lifetime is a pre-deployment test. The third, an agent reusing the human's
root credential (SCOPE-03), is a property of how the deployment provisions
credentials, not of any single token, and so is runtime. Least-privilege frameworks
enforce the scope constraint at runtime with monotonic
confinement~\cite{agentauth2025progent}, and measurements show over-privilege is
common in practice~\cite{agentauth2025privmgmtmcp}; our point is that the carried
scope and lifetime are checkable before the agent ever runs.

\noindent\textbf{Agent identity (AID-01/02).} An agent needs a verifiable identity
distinct from the human's so its actions can be attributed and its access revoked.
The exploit: treating a session token as authentication (AID-01), or giving the agent
no workload identity at all (AID-02), leaves no principal to revoke or audit. The
defense is a deployed workload-identity substrate (WIMSE, SPIFFE/SPIRE) that issues
each agent instance a short-lived, verifiable
credential~\cite{agentauth2026aiidentity,agentauth2026ehv,agentauth2025saga}. Both
are runtime: with no identity layer deployed there is nothing in a presented token to
check against, which is why they sit on the runtime side of the line.

\noindent\textbf{Confused deputy (CD-01/02).} An OAuth proxy sits between many clients
and a resource. The exploit: if the proxy reuses one consent across clients, or
matches \texttt{redirect\_uri} loosely, one client inherits another's authorization
(CD-01); and if it exposes dynamic client registration without vetting (RFC~7591), an
attacker registers a client the proxy then trusts (CD-02). These are properties of a
live consent flow and of how the proxy matches redirects at request time, so they
cannot be decided from a token or a static manifest. Protocol-level analyses of MCP
surface exactly these trust-propagation
gaps~\cite{agentauth2026breakingprotocol,agentauth2025gluecode}, yet across the
literature the confused deputy appears almost only as a named taxonomy entry, not an
enforced property (\S\ref{sec:related}).

\noindent\textbf{Agency (AGT-01/02).} The agency threats are where authorization
meets the model. The exploit: untrusted input drives the agent to escalate privilege
(AGT-01) or to call a tool whose description has been poisoned (AGT-02), so the agent
itself issues the malicious request with its own legitimate token. No token check
helps, because the token is valid and the request is the agent's own; the decision
depends on the executing context. The defenses are runtime: per-action authorization
or a human in the loop so an injection cannot escalate on its
own~\cite{agentauth2025progent}, and vetting of tool descriptions before they reach
the model~\cite{agentauth2026trustdesc,agentauth2026unicodetag}. This category is the
most densely studied in the literature (\S\ref{sec:related}), which is part of our
point: the field has concentrated on the model-level attacks and under-formalized the
token-level guarantees.

\noindent\textbf{A worked scenario.} Consider a travel-planning agent that books
flights through an MCP server, which in turn calls an airline API. The user
authenticates once; the agent then receives an access token and acts autonomously.
Three failures from Table~\ref{tab:tax} compose. If the booking server forwards the
user's token to the airline API rather than exchanging it (MCP-01), and the airline
shares the issuer, a malicious or compromised airline endpoint replays the user's
token against any resource that issuer serves. If the token carries scope
\texttt{travel:*} rather than \texttt{flights:book} (SCOPE-01), a prompt injected
through a hotel-review page the agent reads can drive it to cancel bookings or read
the user's itinerary history, because the token already permits it. And if the
delegated token carries no \texttt{act} claim (DEL-01), neither the airline nor an
auditor can later establish that the agent, rather than the user directly, issued the
request. Each of these three is visible in the token and the server's declared
contract \emph{before the agent runs}: a defender mints a token with the wrong
audience, an over-broad scope, and a missing \texttt{act}, and confirms the verifier
rejects each. The per-client human consent the airline relied on (CD-01) and the prompt injection
that triggered the misuse (AGT-01) are not visible this way; they need a live human
approver and a running agent. The same scenario therefore straddles the line, and the
line tells the team which part to gate in CI, which part to test against staging, and
which part needs a runtime control.

\section{Systematizing the space}\label{sec:syst}
Our systematization rests on three axes that, together, no prior work applies to
agent authorization (\S\ref{sec:related}).

\noindent\textbf{(1) Pre-deployment testability.} Testability before an agent runs is
not one thing but a short spectrum, and collapsing its halves is a common imprecision
we make explicit. We separate three modes. \emph{Token inspection} decides a property
from a presented token and the issuer's key set alone: whether an \texttt{aud} is
present and names this server, whether \texttt{scope} lies within a declared
allow-list, whether \texttt{exp} is set and within a cap, and whether a required
\texttt{act} claim is present and well-formed. \emph{Conformance testing} mints the
adversarial token an attacker would and confirms that the \emph{deployed verifier}
rejects it; it runs before the agent is exercised but, unlike token inspection, needs
the server's enforcement path, because whether a server honors an \texttt{aud} or
refuses passthrough is a property of its implementation, not of the token.
\emph{Runtime monitoring} (\textsf{H}) needs a live human approver or the model
itself. These three modes, not a binary, are what Table~\ref{tab:tax}'s Test column
records, and by them the fifteen threats sort \textbf{6/6/3}. Six are token-static
(\textsf{T}): audience binding (MCP-01/02), the delegation \emph{hygiene} checks
(DEL-01/02, in the bounded sense of \S\ref{sec:tax}), and scope and expiry
(SCOPE-01/02). Six more are server-dynamic (\textsf{S}): whether the issuer enforces
\texttt{may\_act} or issues delegation rather than impersonation (DEL-03/04), whether
the agent is provisioned a distinct credential (SCOPE-03), whether a workload-identity
substrate attests each agent (AID-01/02), and whether dynamic client registration is
left open (CD-02). Each of the six is testable against a staging server or issuer with
no human and no model in the loop; a binary static/runtime split understates this by
filing them all under ``runtime.'' Only three are irreducibly \textsf{H}: per-client
human consent (CD-01) and the two agency threats, prompt-injection escalation and tool
poisoning (AGT-01/02), which turn on the model's behavior under untrusted input. The
result is sharper than a two-way line suggests: twelve of fifteen threats are testable
before production, six from a token and six against a staging server, and only three
need a human or the model.

\noindent\textbf{(2) Normative provenance.} Each threat maps to the clause it
violates: audience binding to RFC~8707 and the MCP passthrough prohibition;
delegation to RFC~8693; scope and expiry to RFC~9700 and RFC~8725; registration to
RFC~7591; confused deputy to CWE-441; agency to the OWASP LLM and Agentic
Top-10s~\cite{owaspllm,owaspagentic}. This lets a practitioner trace a defense to the
authority that mandates it, a mapping the surveyed literature does not provide.

\noindent\textbf{(3) The enforcement spectrum.} The defenses in the
literature scatter across a spectrum they never name: \emph{static} enforcement at the
trust boundary (control/data-flow extraction~\cite{agentauth2025camel}, capability
attestation, contract checks on tokens and key sets) versus \emph{runtime}
enforcement during execution (monotonic privilege confinement~\cite{agentauth2025progent},
pre-action authorization~\cite{agentauth2026preaction}, mandatory access
control~\cite{agentauth2026mac}, out-of-band policy~\cite{agentauth2026angels}).
Drawing the line explicitly reframes several ``runtime-only'' defenses as
compensating controls for threats a pre-deployment contract check could have
excluded: if audience binding and scope are verified statically, a runtime
controller has strictly less to contain. Figure~\ref{fig:line} places the fifteen
threats on the three-way spectrum this yields, \textsf{T}/\textsf{S}/\textsf{H}.

\begin{figure*}[t]
\centering
\begin{tikzpicture}[font=\footnotesize]
\fill[black!12] (0,-.75) rectangle (5.2,1.2);
\fill[black!6]  (5.3,-.75) rectangle (10.5,1.2);
\draw[very thick] (5.25,1.25)--(5.25,-.8);
\draw[very thick] (10.55,1.25)--(10.55,-.8);
\node[anchor=north west,align=left] at (0.1,1.17){\textbf{\textsf{T} token}: a token, key set, contract};
\node[anchor=north west,align=left] at (5.4,1.17){\textbf{\textsf{S} staging server}: no human, no model};
\node[anchor=north west,align=left] at (10.65,1.17){\textbf{\textsf{H} human / model}};
\node[anchor=north west,align=left] at (0.1,.6){MCP-01, MCP-02\\DEL-01, DEL-02\\SCOPE-01, SCOPE-02};
\node[anchor=north west,align=left] at (5.4,.6){DEL-03, DEL-04\\SCOPE-03, CD-02\\AID-01, AID-02};
\node[anchor=north west,align=left] at (10.65,.6){CD-01\\AGT-01\\AGT-02};
\node[anchor=north west,align=left] at (0.1,-.35){\emph{6: testable from a token}};
\node[anchor=north west,align=left] at (5.4,-.35){\emph{6: testable against staging}};
\node[anchor=north west,align=left] at (10.65,-.35){\emph{3: need a human/model}};
\end{tikzpicture}
\caption{The testability spectrum. Twelve of fifteen threats are testable before
production, six from a token alone (\textsf{T}) and six against a staging server or
issuer with no human or model in the loop (\textsf{S}); only three (\textsf{H}) need a
live human approver or the model. This three-way cut, sharper than a binary
static/runtime line, says what a pre-deployment tool, and separately a tightened
standard, can and cannot guarantee.}
\label{fig:line}
\end{figure*}

\noindent\textbf{An instrument for the token-inspection subset.} We release an
open-source analyzer that performs the token-inspection checks of mode one: it decodes
a presented token against a policy and reports cited findings on audience binding,
scope, expiry, and the delegation hygiene property (a required \texttt{act} claim
present and well-formed). It is explicit about the boundary the three modes draw. It
\emph{verifies} the properties a token and key set settle; it only \emph{decodes and
reports} the rest, because, exercising no server, it cannot attest that a deployment
refuses passthrough or enforces \texttt{may\_act}. That attestation is conformance
testing, for which MCP security benchmarks and harnesses are the right
instrument~\cite{agentauth2025mcpsecbench,agentauth2025mcpxkit} and which our
token-level checks complement rather than replace. The analyzer is stateless and makes
no network calls, so it runs in CI or beside the agent. It covers the six token-static
threats; the six server-dynamic threats need a conformance harness against a staging
server, and the three human-or-model threats a runtime control. Naming that boundary
is itself the result: it states what a pre-deployment token check, and separately a
conformance suite and a tightened standard, can and cannot guarantee.

\section{Evaluation: do real verifiers reject the tokens?}\label{sec:eval}
The server-dynamic claim is testable, so we tested it. We stood up nine independent
verifiers: eight widely-used libraries spanning three ecosystems, Go (golang-jwt with
keyfunc, lestrrat-go/jwx, go-jose), Node (jsonwebtoken with jwks-rsa, panva/jose), and
Python (PyJWT, Authlib, python-jose), plus the Envoy \texttt{jwt\_authn} proxy that
fronts production service meshes. Each was configured to its standard access-token
validation: verify the RS256 signature against a test issuer's JWKS, require the issuer
and the audience, and require expiry where the verifier offers an option for it. A
correct resource server must reject every adversarial token. We minted the token-static
corpus of \S\ref{sec:syst} bound to each verifier, fired it as a bearer token, and
recorded acceptances. The generator and the verifiers share no code, so this is an
independent test, not the generator checking its own analyzer.

Every verifier accepted the valid control, confirming it trusts the issuer. \textbf{Six
of the nine then accepted a token they should have rejected} (Table~\ref{tab:eval}).
Four accepted a token carrying no \texttt{exp}, among them the Envoy proxy, not only
libraries: these verifiers offer no way to \emph{require} expiry in their standard
configuration, so a token that never expires passes, against RFC~8725. One, python-jose,
accepted a token with no bound audience even though the expected audience was supplied,
which is the unbound-audience passthrough of \S\ref{sec:tax}; we reproduced it outside
the harness. One accepted a structurally malformed token. Every other attack, algorithm
confusion and the \texttt{alg:none} family, tampered and empty signatures,
\texttt{kid}/\texttt{jku}/\texttt{x5u} injection, wrong issuer or audience, expired and
not-yet-valid, was rejected by all nine. These are defaults, not exotic bugs, which is
the systematization's point made concrete: the token-level weaknesses are the common
case, and a widely-deployed proxy is among those that miss one. The run is also the
mode-\textsf{S} conformance test of \S\ref{sec:syst} in practice, separating the
verifiers that enforce the token-static invariants from those that do not.

\begin{table}[t]
\centering\footnotesize
\caption{Nine reference verifiers against the token-static corpus, standard validation
config. A token accepted is a token a correct resource server must reject. Four accept
a token with no \texttt{exp} (SCOPE-02), including the Envoy proxy; one accepts an
unbound audience (MCP-01/02 passthrough); one accepts a malformed token.}
\label{tab:eval}
\begin{tabular}{@{}llc@{}}
\toprule
Verifier & Kind & Accepted \\
\midrule
golang-jwt + keyfunc & Go lib & --- \\
lestrrat-go/jwx      & Go lib & malformed token \\
go-jose/go-jose      & Go lib & no-\texttt{exp} token \\
jsonwebtoken + jwks-rsa & Node lib & no-\texttt{exp} token \\
jose (panva)         & Node lib & no-\texttt{exp} token \\
PyJWT                & Python lib & --- \\
Authlib              & Python lib & --- \\
python-jose          & Python lib & unbound audience \\
Envoy \texttt{jwt\_authn} & proxy & no-\texttt{exp} token \\
\bottomrule
\end{tabular}
\end{table}

\section{Related work}\label{sec:related}
We surveyed 44 research works on authorization, delegation, identity, and token
security for agents and MCP, published 2023--2026 and gathered by keyword and
citation search over arXiv and the OAuth, MCP, OWASP, and CWE corpora; all arXiv
identifiers were verified to exist. We coded each work for which of the seven threat
categories of Table~\ref{tab:matrix} it treats and at what depth (central
contribution, partial, or absent), which yields the coverage scoring used below. The
fifteen threats are the authorization failures these works name that are expressible
as a property of a token, a key set, or a declared contract; we rate a threat
\emph{critical} when it directly grants an injected agent authority beyond its task or
defeats attribution, \emph{high} when it materially widens exposure, and \emph{medium}
when it is a precondition that enables the others rather than causing direct loss,
following the severity bands of the OWASP LLM and Agentic risk
ratings~\cite{owaspllm,owaspagentic}. The works divide into \emph{threat taxonomies}
and \emph{point mechanisms}, and neither systematizes authorization by our three axes.

\noindent\textbf{Taxonomies and surveys.} MCP security surveys enumerate attack
surfaces across the server lifecycle~\cite{agentauth2025mcplandscape}, and recent
systematizations consolidate attacks and defenses for the
ecosystem~\cite{agentauth2025sokmcp} and for multi-agent systems end to
end~\cite{agentauth2026sokmultiagent}. Cross-protocol surveys map the security and
interoperability of MCP, A2A, ACP, and ANP and compare their
guarantees~\cite{agentauth2025commsurvey,agentauth2025interopsurvey,agentauth2025protocolcompare}.
Threat taxonomies name authorization-adjacent
risks explicitly, including the confused deputy and over-broad
consent~\cite{agentauth2026mcp38}, and a defense-placement taxonomy classifies where
in the architecture a control sits~\cite{agentauth2026mcpdpt}. Protocol-level
analyses dissect the MCP specification itself, finding missing capability attestation
and implicit trust propagation across chained
servers~\cite{agentauth2026breakingprotocol,agentauth2026mcpdesign}. These organize by
attacker type, lifecycle phase, defense placement, or ecosystem component, and none
sorts threats by pre-deployment testability or binds each to the clause it violates.

\noindent\textbf{Delegation and identity mechanisms.} A line of work makes delegation
verifiable: authenticated-delegation frameworks give an agent scoped, auditable
authority to act for a user~\cite{agentauth2025authdelegation}; Agentic JWT binds each
agent action to a user intent in the token~\cite{agentauth2025agenticjwt}; AIP carries
delegation chains across MCP and A2A boundaries~\cite{agentauth2026aip}; bounded-delegation and stateful-effect schemes
keep a sub-agent within what was delegated~\cite{agentauth2026boundedagents,agentauth2026aidguard};
and secret-use delegation and residual authorization treat scoped secret handling and
revocation~\cite{agentauth2026sudp,agentauth2026residualauth}. Identity work gives
agents a verifiable, distinct principal through WIMSE and
SPIFFE/SPIRE~\cite{agentauth2026aiidentity,agentauth2026ehv,agentauth2026ldp}, and
full governance architectures combine registration, identity, and access
tokens~\cite{agentauth2025saga}. These are the mechanisms our delegation and identity
categories draw on, each addressing a slice rather than the authorization layer as a
whole.

\noindent\textbf{Privilege control and model-level defenses.} Runtime privilege
frameworks enforce least privilege with symbolic rules and monotonic
confinement~\cite{agentauth2025progent}, move the authorization decision before each
tool call~\cite{agentauth2026preaction}, apply mandatory access control to agent tool
use~\cite{agentauth2026mac}, isolate tools in separate
environments~\cite{agentauth2024isolategpt}, and enforce policy out-of-band at a
trusted boundary~\cite{agentauth2026angels}; measurements show over-privilege is
widespread in practice~\cite{agentauth2025privmgmtmcp}, privilege escalation through
the agent's execution context has its own analysis~\cite{agentauth2026contextroot},
and runtime authorization governs the resources an agent acquires
mid-task~\cite{agentauth2026runtimeauthz}. For the model-level threats,
by-design defenses extract control and data flow so untrusted input cannot alter
program flow~\cite{agentauth2025camel}, benchmarks quantify injection
success~\cite{agentauth2024agentdojo}, and tool-poisoning work measures and defends
the description surface~\cite{agentauth2025mcptox,agentauth2026trustdesc,agentauth2026noboxmcp,agentauth2026unicodetag},
with indirect prompt injection the seminal threat~\cite{agentauth2023indirectpi}, and
recent taxonomies chart the prompt-injection threat space for LLM agents and for
agentic coding assistants
specifically~\cite{agentauth2026pitaxonomy,agentauth2026codingassistants}. Each
is one point on the static/runtime line; none draws the line, and the densest coverage
is on the model-level attacks rather than the token-level guarantees.

\noindent\textbf{The coverage gap, quantified.} Scoring the representative works on
seven threat categories (Table~\ref{tab:matrix}) makes the gap numeric rather than
impressionistic: of the twenty-one works scored, prompt injection is a central
contribution in fourteen, but audience binding is central in only three (Agentic
JWT~\cite{agentauth2025agenticjwt}, aiAuthZ~\cite{agentauth2026aiauthz}, and one
enterprise framework~\cite{agentauth2025enterprisemcp}) and confused-deputy / dynamic
registration in only two~\cite{agentauth2026mcp38,agentauth2025enterprisemcp}, where it
figures mostly as a named taxonomy entry rather than a verified or enforced property.
The field has thoroughly characterized the model-level attacks and under-formalized the
token- and consent-level authorization guarantees, the seam this SoK occupies. This
gap is a claim about the \emph{literature}, evidenced by the survey's own coding;
whether the corresponding weaknesses are also prevalent in deployed servers is a
separate, empirical question we scope below.

\begin{table*}[t]
\centering\footnotesize
\caption{Coverage of the seven agent-authorization threat categories across
representative works. \yes: a central contribution; $\circ$: partially addressed;
blank: not addressed. AudBind, audience binding and passthrough; OBO, on-behalf-of
delegation; Scope, scope minimization; ConfDep, confused deputy and dynamic
registration; WID, workload/agent identity; PI, prompt-injection escalation; Tool,
tool poisoning. Two columns, AudBind and ConfDep, are thin.}
\label{tab:matrix}
\begin{tabular}{@{}lccccccc@{}}
\toprule
Work & AudBind & OBO & Scope & ConfDep & WID & PI & Tool \\
\midrule
\multicolumn{8}{@{}l}{\emph{Taxonomies and surveys}}\\
MCP Landscape~\cite{agentauth2025mcplandscape}      & $\circ$ & $\circ$ & $\circ$ & $\circ$ & $\circ$ & \yes & \yes \\
SoK MCP~\cite{agentauth2025sokmcp}                  & $\circ$ & $\circ$ & $\circ$ & $\circ$ & $\circ$ & \yes & \yes \\
MCP-38~\cite{agentauth2026mcp38}                    & $\circ$ & $\circ$ & $\circ$ & \yes    & $\circ$ & \yes & \yes \\
Enterprise MCP~\cite{agentauth2025enterprisemcp}    & \yes    & \yes    & \yes    & \yes    & $\circ$ & \yes & \yes \\
Breaking the Protocol~\cite{agentauth2026breakingprotocol} & $\circ$ & & & $\circ$ & $\circ$ & \yes & \yes \\
\midrule
\multicolumn{8}{@{}l}{\emph{Delegation and identity mechanisms}}\\
Agentic JWT~\cite{agentauth2025agenticjwt}          & \yes    & \yes    & \yes    & $\circ$ & \yes    & $\circ$ & \\
AIP~\cite{agentauth2026aip}                         & $\circ$ & \yes    & $\circ$ & $\circ$ & \yes    & & \\
Bounded Agents~\cite{agentauth2026boundedagents}    & $\circ$ & \yes    & \yes    & $\circ$ & $\circ$ & $\circ$ & \\
AID-Guard~\cite{agentauth2026aidguard}              & $\circ$ & \yes    & \yes    & $\circ$ & $\circ$ & $\circ$ & \\
SAGA~\cite{agentauth2025saga}                       & $\circ$ & \yes    & \yes    & $\circ$ & \yes    & $\circ$ & \\
AI Identity~\cite{agentauth2026aiidentity}          & $\circ$ & \yes    & $\circ$ & $\circ$ & \yes    & & \\
aiAuthZ~\cite{agentauth2026aiauthz}                 & \yes    & $\circ$ & $\circ$ & $\circ$ & \yes    & & \\
\midrule
\multicolumn{8}{@{}l}{\emph{Privilege / runtime control}}\\
Progent~\cite{agentauth2025progent}                 & & & \yes & & & \yes & $\circ$ \\
Pre-Action Authz~\cite{agentauth2026preaction}      & $\circ$ & $\circ$ & \yes & $\circ$ & $\circ$ & \yes & \\
MAC framework~\cite{agentauth2026mac}               & & $\circ$ & \yes & $\circ$ & & \yes & $\circ$ \\
IsolateGPT~\cite{agentauth2024isolategpt}           & & & \yes & $\circ$ & $\circ$ & \yes & $\circ$ \\
If Agents Were Angels~\cite{agentauth2026angels}    & $\circ$ & $\circ$ & \yes & $\circ$ & $\circ$ & \yes & $\circ$ \\
\midrule
\multicolumn{8}{@{}l}{\emph{Model-level attacks and defenses}}\\
Indirect PI~\cite{agentauth2023indirectpi}          & & & & & & \yes & $\circ$ \\
CaMeL~\cite{agentauth2025camel}                     & & & \yes & $\circ$ & & \yes & $\circ$ \\
MCPTox~\cite{agentauth2025mcptox}                   & & & $\circ$ & & & \yes & \yes \\
No-Box MCP~\cite{agentauth2026noboxmcp}             & & & & & & \yes & \yes \\
\bottomrule
\end{tabular}
\end{table*}

\section{Discussion: what a standard could mandate}\label{sec:disc}
The testability spectrum has a direct standardization reading: the twelve testable
threats are exactly the ones a profile could make mandatory-to-verify. An MCP or OAuth
profile that \emph{requires} resource indicators (RFC~8707), rejects passthrough,
mandates nested \texttt{act} validation, enforces \texttt{may\_act} on exchange, and
caps scope and lifetime would move them from ``deployer's responsibility, usually
missed'' to ``checked at the boundary.'' The split tells a conformance suite how to
check each: the six token-static requirements it adjudicates from a token alone, and
the six server-dynamic ones by firing minted tokens, registration requests, and token
exchanges at a candidate server and recording which it accepts. Table~\ref{tab:conformance}
makes the token-static half concrete, pairing each clause with the token minted to
violate it and the expected rejection, so the profile ships with its own executable
check rather than prose a deployer interprets. Only the three human-or-model threats
fall outside what such a profile and suite can settle.

\begin{table}[t]
\centering\footnotesize
\caption{A conformance profile for the six token-static threats. Each row is a
requirement a profile could mandate and a test a suite runs from a token, a key set,
and a declared contract, with no running agent. The six server-dynamic threats extend
this with tests fired at a candidate server.}
\label{tab:conformance}
\begin{tabular}{@{}llp{3.3cm}@{}}
\toprule
Threat & Clause & Conformance test \\
\midrule
MCP-01/02 & RFC~8707 & Mint a token whose \texttt{aud} names another resource; the
verifier must reject it. \\
DEL-01/02 & RFC~8693 & Present a delegated token with a missing or malformed
\texttt{act} chain; the verifier must reject it. \\
SCOPE-01 & RFC~9700 & Present a token whose scope exceeds the tool's declared need;
the contract check must flag it. \\
SCOPE-02 & RFC~8725 & Present a token past its \texttt{exp} or beyond a capped
lifetime; the verifier must reject it. \\
\bottomrule
\end{tabular}
\end{table}

\noindent\textbf{An agenda for the untested nine.} The nine threats outside token
inspection are not uniformly open; the spectrum sorts them into two agendas. Six are
server-dynamic and want \emph{engineering}: a conformance harness that fires minted
tokens and registration and exchange requests at a candidate server and records what
it accepts would cover whether the issuer enforces \texttt{may\_act} or issues
delegation rather than impersonation (DEL-03/04), whether the agent is provisioned a
distinct credential (SCOPE-03), whether a deployed substrate such as
SPIFFE/SPIRE~\cite{agentauth2026ldp} attests each agent (AID-01/02), and whether
dynamic registration is left open (CD-02). The other three want \emph{mechanism
research}, because no server-side test settles them: per-client human consent (CD-01)
needs an auditable consent record so that which client obtained which grant is
reconstructable rather than trusted, and the agency threats (AGT-01/02) need a trusted
mediation boundary that survives a prompt-injected plan, which the control- and
data-flow work~\cite{agentauth2025camel} approaches but does not yet reduce to a
standard. Separating the engineering agenda from the research one is the practical
payoff of the spectrum: it says which untested threats a conformance suite can retire
and which remain genuinely open.

\section{Threats to validity}\label{sec:validity}
The MCP specification is fast-moving; this systematization is a dated snapshot. The
testability spectrum depends on token and manifest availability; a deployment that
hides relevant state reduces what any check can conclude. The token/server/human cut
is a judgment on the weakest harness per threat; a few threats sit near a boundary
(for example AID-01/02 are server-dynamic only once an identity substrate is deployed
in staging), which we mark rather than hide. Our survey covers
research indexed by 2026-10; most agent-authorization work is recent and not yet
peer-reviewed, which we mark. The coverage gap of \S\ref{sec:related} is a claim about
the research literature; our evaluation (\S\ref{sec:eval}) gives direct evidence that
the token-static weaknesses occur in widely-used verifier libraries, not only on paper.
\gate{A broader prevalence measurement over a corpus of public MCP server deployments,
with a hand-validated stratified sample, reported inter-rater agreement, and coordinated
disclosure for identifiable high-severity findings, remains future work.}

\section*{Ethics considerations}
The systematization uses only public specifications, advisories, and CVEs. The
accompanying analyzer runs on tokens the operator already holds and sends nothing
anywhere; we probe no third-party agents or servers and use no real user tokens. Any
future measurement will characterize population prevalence, not single deployments,
and will follow coordinated disclosure for identifiable high-severity findings. No
personal data or human subjects are involved.

\section*{Open science}
We release the static analyzer, a machine-readable threat taxonomy (each threat with
its invariant and normative sources), and the evaluation harness with its
nine-verifier target set and raw results, under an open licence, so both the
systematization and the evaluation of \S\ref{sec:eval} are reproducible.
\gate{A DOI-pinned artifact will accompany the camera-ready version.}

\section{Conclusion}
Agent authorization is a high-stakes surface the research literature has studied
around rather than through. We systematize it by pre-deployment testability, by
normative provenance, and by a three-way enforcement spectrum; we show that twelve of
fifteen threats are testable before production, six from a token alone and six against
a staging server with no human or model in the loop, and only three need a live human
approver or the model itself; we evaluate the token-static subset against nine
reference verifiers and find six accept a token they should reject, so the weaknesses
are defaults rather than corner cases; and we argue this cut should guide what a
pre-deployment tool checks, what a conformance suite exercises, and what a standard
makes mandatory to verify.

\section*{LLM Usage Statement}
LLMs were used for editorial purposes in this manuscript, and all outputs were
inspected by the authors to ensure accuracy and originality. An LLM-assisted process
additionally supported the literature survey of \S\ref{sec:related}: candidate works
were gathered with model assistance and each was then verified by the authors to
exist, with titles and author lists confirmed against the arXiv API on 2026-10-06;
works that could not be verified were excluded, and one paper that its own authors had
formally withdrawn was identified this way and removed. The authors are responsible
for the selection of works, the coverage judgments in the systematization, and every
claim in the paper. One reproducibility limitation applies: the assisting models are
not openly available, so the survey's search recall cannot be reproduced exactly,
although every cited work is independently verifiable by its arXiv identifier.

\appendix
\section{Worked token-inspection examples}\label{app:examples}
Because the mode-one checks are definitional, their verdicts are deterministic rather
than learned. Table~\ref{tab:examples} gives illustrative adversarial tokens (claims
abridged) and the analyzer's token-inspection verdict against a fixed policy that
requires the audience \texttt{api.pay}, a \texttt{scope} within the allow-list
\{\texttt{files:read}\}, an \texttt{exp} within fifteen minutes, and a well-formed
\texttt{act} claim on any token that asserts delegation. The examples demonstrate what
the analyzer \emph{verifies} and, in the last row, what it does not: a passing
token-inspection verdict does not attest that the deployed verifier enforces the same
constraints, which is conformance testing (\S\ref{sec:syst}). This is a demonstration
of the checks, not an empirical evaluation; measuring how often real deployments carry
such tokens is the prevalence study scoped in \S\ref{sec:validity}.

\begin{table}[h]
\centering\footnotesize
\caption{Illustrative token-inspection verdicts against the fixed policy of
\S\ref{app:examples}. Verdicts are deterministic because the checks are definitional.}
\label{tab:examples}
\begin{tabular}{@{}llp{3.3cm}@{}}
\toprule
ID & Token (abridged) & Verdict \\
\midrule
A1 & \texttt{aud}=api.other & FAIL MCP-02: \texttt{aud} names another resource \\
A2 & \texttt{scope}=admin:* & FAIL SCOPE-01: scope exceeds declared need \\
A3 & no \texttt{exp} & FAIL SCOPE-02: no lifetime bound \\
A4 & \texttt{act} absent, asserts OBO & FAIL DEL-01: delegation hygiene, \texttt{act} missing \\
A5 & all constraints met & PASS inspection; does not attest server enforcement \\
\bottomrule
\end{tabular}
\end{table}
